\chapter{Boolean Analysis --- Cube}
%%% generated by the blueprint pipeline from TCSlib.BooleanAnalysis.GaussianSpace.Bobkov.Cube
\section{Overview}
Coordinate restrictions respect discrete derivatives and split the squared gradient. Euclidean triangle inequalities give expectation and midpoint bounds used in tensorization.

\section{Declarations}

\begin{lemma}[Derivative under last-coordinate restriction]
\lean{BooleanAnalysis.GaussianSpace.derivative_restrictLast}
\leanok
\uses{BooleanAnalysis.BoolCube, BooleanAnalysis.BooleanFunc, BooleanAnalysis.Hypercontractivity.restrictLast, BooleanAnalysis.derivative}
\difficulty{1}
Let $n$ be a natural number, let $f:\{0,1\}^{n+1}\to\mathbb R$, let $i$ be a coordinate
in $\{0,\ldots,n-1\}$, let $b$ be a bit, and let $x\in\{0,1\}^n$. The discrete
derivative in direction $i$ of the restriction of $f$ to last coordinate $b$, evaluated
at $x$, equals the discrete derivative of $f$ in coordinate $i$ at the point obtained by
appending $b$ to $x$.
\end{lemma}

\begin{lemma}[Last derivative at an appended point]
\lean{BooleanAnalysis.GaussianSpace.derivative_last_snoc}
\leanok
\uses{BooleanAnalysis.BoolCube, BooleanAnalysis.BooleanFunc, BooleanAnalysis.Hypercontractivity.diffLast, BooleanAnalysis.Hypercontractivity.restrictLast, BooleanAnalysis.derivative}
\difficulty{1}
Let $n$ be a natural number and let $f:\{0,1\}^{n+1}\to\mathbb R$. For every bit $b$ and
every $x\in\{0,1\}^n$, the discrete derivative of $f$ in the last coordinate at the
point obtained by appending $b$ to $x$ equals the last-coordinate half-difference of $f$
at $x$.
\end{lemma}

\begin{lemma}[Squared derivatives under last-coordinate restriction]
\lean{BooleanAnalysis.GaussianSpace.sum_sq_derivative_snoc}
\leanok
\uses{BooleanAnalysis.BoolCube, BooleanAnalysis.BooleanFunc, BooleanAnalysis.GaussianSpace.derivative_last_snoc, BooleanAnalysis.GaussianSpace.derivative_restrictLast, BooleanAnalysis.Hypercontractivity.diffLast, BooleanAnalysis.Hypercontractivity.restrictLast, BooleanAnalysis.derivative}
\difficulty{3}
Let $n\ge 0$, let $f:\{0,1\}^{n+1}\to\mathbb R$, let $b\in\{0,1\}$, and let
$x\in\{0,1\}^n$. The sum of the squares of all coordinate half-differences of $f$ at
$(x,b)$ equals the sum of the squares of the coordinate half-differences of the
restriction $f_b(z)=f(z,b)$ at $x$, plus the square of the last-coordinate
half-difference of $f$ at $x$:
\[\sum_{i=0}^{n}(D_i f(x,b))^2=\sum_{i=0}^{n-1}(D_i f_b(x))^2+\left(\frac{f(x,0)-f(x,1)}{2}\right)^2.\]
\end{lemma}

\begin{lemma}[Expectation of the last-coordinate half-difference]
\lean{BooleanAnalysis.GaussianSpace.expect_diffLast}
\leanok
\uses{BooleanAnalysis.BooleanFunc, BooleanAnalysis.Hypercontractivity.diffLast, BooleanAnalysis.Hypercontractivity.restrictLast, BooleanAnalysis.expect}
\difficulty{3}
Let $n$ be a natural number and let $f\colon\{0,1\}^{n+1}\to\mathbb R$. For
$b\in\{0,1\}$, write $f_b(x)=f(x,b)$ for $x\in\{0,1\}^n$, and set
$d(x)=\tfrac12(f_0(x)-f_1(x))$. With expectation taken uniformly over the Boolean cube,
$\mathbb E[d]=\tfrac12(\mathbb E[f_0]-\mathbb E[f_1])$.
\end{lemma}

\begin{lemma}[Convexity of the Euclidean norm under expectation]
\lean{BooleanAnalysis.GaussianSpace.sqrt_expect_sq_add_sq_le_expect_sqrt}
\leanok
\uses{BooleanAnalysis.BoolCube, BooleanAnalysis.BooleanFunc, BooleanAnalysis.expect, BooleanAnalysis.expect_eq_fintypeExpect}
\difficulty{4}
For every dimension $n$ and real-valued functions $u,v$ on the Boolean cube
$\{0,1\}^n$, the Euclidean norm of their expectations is at most the expectation of
their pointwise Euclidean norm: $\sqrt{(\E[u])^2+(\E[v])^2}\le\E[\sqrt{u^2+v^2}]$.
\end{lemma}

\begin{lemma}[Midpoint inequality with Euclidean coordinates]
\lean{BooleanAnalysis.GaussianSpace.carried_midpoint}
\leanok
\difficulty{5}
Let $m$ be a natural number, let $u,v\in\mathbb R^m$, and let $A,B,C\in\mathbb R$
satisfy $A\ge0$ and $A\le(B+C)/2$. Then
\[
\sqrt{A^2+\sum_{j=1}^m
u_j^2}\le\frac{\sqrt{B^2+\sum_{j=1}^m(u_j+v_j)^2}+\sqrt{C^2+\sum_{j=1}^m(u_j-v_j)^2}}{2}.
\]
\end{lemma}

\begin{lemma}[Linearity of the discrete derivative]
\lean{BooleanAnalysis.GaussianSpace.derivative_linear_combination}
\leanok
\uses{BooleanAnalysis.BooleanFunc, BooleanAnalysis.derivative, BooleanAnalysis.derivative_add, BooleanAnalysis.derivative_smul}
\difficulty{3}
Let $n$ be a natural number, let $i$ be a coordinate in $\{0,\ldots,n-1\}$, let
$f,g:\{0,1\}^n\to\mathbb R$, and let $c,d\in\mathbb R$. The discrete derivative in
direction $i$ of the function $x\mapsto c f(x)+d g(x)$ is the function $x\mapsto c(D_i
f)(x)+d(D_i g)(x)$.
\end{lemma}
